% !PRINTED 542
\weberpage{542}

\begin{center}
\S.\,164.\\
\gesperrt{Resolvents of Lagrange.}
\end{center}

The method of solving cyclic equations with which we now wish to become acquainted is applicable uniformly to prime degrees and to composite degrees. For this purpose one makes use of certain expressions, known under the name of the \gesperrt{resolvents of Lagrange}\footnote{\gesperrt{Lagrange}, Réflexions etc., see p. 508. Earlier we understood by resolvents solving equations; here they are solving functions.}, which are of great use in all investigations concerning the algebraic solution of equations.

Let $F(x) = 0$ be an equation with the roots $\alpha, \alpha_1, \alpha_2 \ldots \alpha_{m-1}$. We denote by $\varepsilon$ any $m$th root of unity and introduce the notation
\weq{1}{(\varepsilon, \alpha) = \alpha + \varepsilon\alpha_1 + \varepsilon^2\alpha_2 + \cdots + \varepsilon^{m-1}\alpha_{m-1}.}
The sums so defined are those which one calls the \gesperrt{Lagrange resolvents}. When these functions are known for all $m$th roots of unity $\varepsilon$, the equation itself is also solved, for by \S.\,133 (6) we have
\weq{2}{\overset{\varepsilon}{\Sigma}\,\varepsilon^k = m \ \text{or} = 0,}
according as $k$ is divisible by $m$ or not, and from this it follows
\weq{3}{m\alpha = \overset{\varepsilon}{\Sigma}(\varepsilon, \alpha),}
in which the sums extend over all $m$th roots of unity $\varepsilon$. One can also express the other roots in the same way:
\weq{4}{m\alpha_k = \overset{\varepsilon}{\Sigma}\varepsilon^{-k}(\varepsilon, \alpha),}
so that in fact everything is reduced to the knowledge of $\varepsilon$ and of the functions $(\varepsilon, \alpha)$.

The sums (2) and (3) may also be represented in a somewhat different manner, since all the $m$th roots of unity are powers of a primitive one among them. If, then, one understands by $\varepsilon$ a fixed primitive $m$th root of unity, and by $\lambda$ an index which runs through a complete system of residues with respect to the modulus $m$, say the
%CONT>>
